\begin{table}[p]
\centering
\begin{threeparttable}
\caption{Baseline and correction components of the five S\&P 500 forecasts.}
\label{tab:sp500-baseline-decomposition}
\small
\setlength{\tabcolsep}{4pt}
\begin{tabular*}{\textwidth}{@{\extracolsep{\fill}}llrrr@{}}
\toprule
Origin & Model & $B\times10^4$ & $A$ & $\widehat V_{t,1}\times10^4$ \\
\midrule
@@ROW_0@@
@@ROW_1@@
@@ROW_2@@
\addlinespace
@@ROW_3@@
@@ROW_4@@
@@ROW_5@@
\addlinespace
@@ROW_6@@
@@ROW_7@@
@@ROW_8@@
\addlinespace
@@ROW_9@@
@@ROW_10@@
@@ROW_11@@
\addlinespace
@@ROW_12@@
@@ROW_13@@
@@ROW_14@@
\bottomrule
\end{tabular*}
\begin{tablenotes}[flushleft]
\footnotesize
\item Notes: All origins are in 2020. Components are calculated from the original predictive draws and corresponding posterior parameters. The identity $\widehat V_{t,1}=BA$ holds before rounding. Baseline and correction parameters are jointly estimated, so this decomposition does not identify the causal effect of either component.
\end{tablenotes}
\end{threeparttable}
\end{table}
